\subsection{Primitives for Guardrail Specification}\label{sec:gdl}

We present the \emph{Guardrails Description Language} (GDL), a declarative language for specifying performance guardrails over kernel behavior.
As discussed in \Cref{subsec:semantics}, a guardrail is given by: (i) \emph{what condition} must hold on the system state it inspects (the monitorable property $\property{G}$), (ii) \emph{what corrective action} is applied when that condition fails ($\action{G}$), and (iii) \emph{which policies} the guardrail applies to ($\policies_G$).
Along with the guardrail specification, the system also needs to know the local guardrail properties $\varphi$ and the global \perfsafety property $\Phi$ that the guardrail is designed to enforce.
Guardrail developers (OS operators or OS policy developers) are responsible for providing all these pieces of information.
Next, we describe the primitives \lang introduces to make this specification intuitive.

A crucial aspect of guardrail specification is determining \emph{when} to evaluate the guardrail.
In principle, any point in the execution sequence of state transitions, decisions, and events can be used for evaluation.
However, there is a cost associated with evaluating the run-time performance property -- of reading the relevant state, computing the condition, and then, taking the action (if needed).
Therefore, the question is {\it when should the system incur the cost of evaluating a guardrail?}
In \lang, users can control this overhead by specifying explicit {\em triggers}, which determine \emph{when} the system should check the guardrail property.

Concretely, a guardrail in \lang is declared using six keywords corresponding \aditya{hanging sentence?}: {\sf for}, {\sf when}, {\sf check}, {\sf action}, {\sf ensures}, and {\sf assert}.
The keyword {\sf for} names the set of policies for which the guardrail is relevant.
{\sf when} specifies the \emph{trigger}: the point in the execution at which the guardrail is evaluated.
{\sf check} encodes the performance condition.
{\sf action} prescribes the corrective action taken when the condition fails.
{\sf ensures} specifies the local guardrail properties $\varphi$ that must hold.
Finally, {\sf assert} specifies the global \perfsafety property $\Phi$ that the guardrail is designed to enforce.

\Cref{tab:gdl-primitives} summarizes the key primitives available in \lang. Below, we describe these primitives in more detail and illustrate their use with examples of guardrails  in \lang for the running examples of CBMM (\Cref{fig:intro-example}) and FLUX (\Cref{fig:flux-example}).

\begin{table}[t]
\centering
\small
\renewcommand{\arraystretch}{1} % or 1.3–1.5 to taste
\setlength{\tabcolsep}{4pt}
\begin{tabular}{c l l}
\toprule
%\textbf{Category} 
& \textbf{Primitive} & \textbf{Description} \\
\midrule
%\multirow{2}{*}{Trigger}
\multicolumn{3}{l}{\bf \textit{Trigger Functions}} \\
  & $\textsf{periodic}(d)$ & Evaluate every $d$ seconds \\
  & $\textsf{on\_call}(f)$ & Evaluate on each invocation of $f$ \\ %\textsf{KFunc} $f$ \\
\midrule
%\multirow{8}{1.4cm}{Inspection Objects}
%\multirow{8}{1.4cm}{Signals}
\multicolumn{3}{l}{\bf \textit{Inspection Objects and Observables}} \\
  % & $\textsf{KObject}$ & Kernel object (page, process, \ldots) \\
  % & $Set[\textsf{KObject}]$ & Collections of kernel objects \\
  & $\mathsf k$ & Kernel object variable \\
  & $\textsf{attribute}(x)$ & Attribute of a \textsf{KObject} or set \\
  & $\textsf{exec\_time}(f)$ & Wall clock execution time of $f$ \\ % \textsf{KFunc} $f$ \\
  & $\textsf{decision}(f)$ & Decision value returned by $f$\\ % \textsf{KFunc} $f$ \\
  & $\textsf{arg}(f, \mathit{name})$ & Arg. value at an invocation of $f$ \\
  %\textsf{KFunc} $f$ invocation \\
  % & $\textsf{read}(\mathit{feat})$ & Read a feature from the in-kernel store \\
  & $\textsf{read}(x)$ & Read feature x from in-kernel store \\
  & \makecell[l]{$\textsf{avg}$/$\textsf{min}$/$\textsf{max}$/$\textsf{sum}(\cdot)$} & Aggregate an observation \\
\midrule
%\multirow{2}{*}{Action}
\multicolumn{3}{l}{\bf \textit{Restorative Actions}}\\
  & \textsf{Policy}($\pi$) & Switch to policy $\pi$ \\
  & $\textsf{Set}(\mathit{flag},\,\mathit{val})$ & Update a control flag or parameter \\
  & $\textsf{Report}$ & Log the violation for offline diagnosis \\
\midrule
\multicolumn{3}{l}{\bf \textit{Domains}}\\
  & $\textsf k \in \textsf{KObject}$ & Kernel Objects (e.g., Pages, processes, \ldots) \\
  & $K \in Set[\textsf{KObject}]$ & Sets of Kernel Objects \\
  & $f \in \textsf{KFunc}$ & Kernel Function \\
  & $\pi \in \Pi$ & OS Policy name \\
\bottomrule
\end{tabular}
\caption{Key primitives in the Guardrails Description Language.}
\label{tab:gdl-primitives}
\end{table}

\begin{dslbox}[title={Guardrail~\theguardcount: Example guardrail for CBMM}, glabel=lst:benefit-correction]
for CBMM
when periodic(2s)
check $\forall P_i,P_j \in Set[{\sf Page}]$ (
      accesses($P_i$) > accesses($P_j$)
      $\implies$read(Benefit[$P_i$]) > read(Benefit[$P_j$])
  )
action $\forall P_i\in Set[{\sf Page}]$ Set$(\Delta{\sf benefit}[P_i], expr)$
\end{dslbox}

\paragraph{Triggers.}
\lang supports two trigger types that cover the two most common classes of performance properties.
The first class of properties specifies continuous conditions that must always hold, such as CPU utilization by kernel daemons remaining below a threshold, or memory compaction overhead staying bounded.
Such properties are not tied to specific state transitions and are best checked at regular intervals.
\lang captures this with the {\sf periodic}$(d)$ trigger, which fires every $d$ seconds.
For example, the guardrail for evaluating the accuracy of estimated benefits in CBMM (see \Cref{lst:benefit-correction}) evaluates the property every 2 seconds.

The second class of properties evaluates the correctness or efficacy of a particular policy \emph{decision}: e.g., whether a flow-size prediction used features that are in-distribution for the model, or whether a page promotion moved a cold page into the hot tier.
For such properties, the natural evaluation point is the moment the decision is produced---at the invocation of the relevant kernel function.
\lang captures this with the {\sf on\_call}$(f)$ trigger, which fires on each call to \textsf{KFunc} $f$.
For example, say the predictor from \Cref{fig:flux-example} is trained on only short flows, then the guardrail fires on each call to the {\tt first\_call} function that classifies the flow (see \Cref{lst:in-dist-short}).
The guardrail then checks whether the input features lie within the training distribution of the predictor.
\begin{dslbox}[title={Guardrail~\theguardcount: Example guardrail to check in-distribution inputs for a neural network trained on short flows}, glabel=lst:in-dist-short]
for FLUXShortPredictor
when on_call(first_call)
check $agg\_net \sim$ read(ShortDist[agg_net]) $\wedge$
      $cpu \sim$ read(ShortDist[cpu]) $\wedge\ \dots$
action Set(flux_policy, FLUXLongPredictor)
ensures flux_policy_is_safe().
assert flux_policy_is_safe().
\end{dslbox}

There may be properties tied to higher-level events, where the state transition of interest may be missed by periodic sampling and cannot be captured cleanly by any single function invocation.
Such cases can be handled by encoding the event as a predicate over observations and evaluating the guardrail periodically at a sufficiently fine granularity.
Concretely, if a guardrail is intended to enforce $\property{}$ each time an event $E$ occurs, one can equivalently specify it with a periodic trigger and the condition $E\implies\property{}$; under a sufficiently small period, $\property{}$ will be evaluated exactly when $E$ holds.

\paragraph{Inspection objects and observations.}
As discussed in \Cref{subsec:property-principles}, guardrails properties must use OS statistics available at runtime, including in-kernel statistics, occurrence of various events, kernel policy decisions, etc.
\lang exposes these through a small set of \emph{inspection objcets} (see \Cref{tab:gdl-primitives}).
These include kernel objects (\textsf{KObject}), such as pages or processes; sets of kernel objects ($Set[\textsf{KObject}]$), which allow a guardrail to reason about collections as first-class entities (e.g., whether a property holds for all memory regions of 256 contiguous pages); and instrumented kernel functions (\textsf{KFunc}), which represent decision points or state transformations.
The choice of these inspection objects is deliberate: guardrails targeted by \lang ultimately reason about kernel objects, their collections, and the function invocations that manipulate them.
By centering the language around these domains, \lang remains narrow while still covering a broad space of runtime performance conditions. \aditya{this sounds unconvincing and vague. can we strengthen the argument for inspection objects? "centering" on "domains" is vague.}

Over these domains, \aditya{what are domains?} operators define \emph{observations} that extract the relevant runtime information.
This includes attributes of \textsf{KObject}s (e.g., {\tt accesses} for a page in \Cref{lst:benefit-correction}), and similarly, set-level attributes of $Set[\ldots]$s (e.g., number of objects in the set).
One can also use observations over decisions and execution times of instrumented functions using the primitives {\tt decision} and {\tt exec\_time} in \Cref{tab:gdl-primitives}.
Any custom metric not representable through \textsf{KObject}s or \textsf{KFunc}s can be exposed via the {\sf read}(\textit{feat}) operation (see \Cref{tab:gdl-primitives}), which reads a pre-populated entry from an in-kernel store.
Note that guardrails are not responsible for collecting these features; the user must implement and maintain the relevant collection mechanisms. 
Also, the correctness of feature collection is a burden of the user, not the guardrails framework.

\lang also supports aggregation ({\sf avg}, {\sf min}, {\sf max}, {\sf sum}) of observations and quantification ($\forall$, $\exists$) over a particular type of inspection objects.
This separation between inspection objects and observations is important: the former determines \emph{what objects} a guardrail looks at, while the latter determines \emph{what it measures} about them.
Operators write the performance condition as predicates over these observations.
For instance, the CBMM guardrail in \Cref{lst:benefit-correction} evaluates whether the accesses between any two sets of pages holds the same relationship as the estimated benefits. This condition ensures that the benefits hold the same relationship as the accesses and are thus, accurate. \aditya{I had a really hard time understanding this subsubsection. also it breaks the flow - can we move it to the end or elsewhere?}
% Thus, although the semantics of guardrails are defined over execution traces, \lang does not expose raw traces directly; instead, it provides a structured interface for expressing performance conditions concisely.

\paragraph{Actions.}
A guardrail must also determine how the system should respond when a property is violated.
While the space of possible corrective actions is broad -- different policies and different classes of violations may require different interventions -- in \lang, we expose only a narrow interface supporting two types of actions.
The first is a non-intrusive action, {\sf Report}, which records whenever a violation occurs for offline diagnosis; such reports can help in designing new policies altogether \ds{Add forward pointer to an example in memory tiering}.
The second is a corrective action, {\sf Set}$(flag, val)$, which modifies a control flag or parameter to alter subsequent system behavior.
Such updates may tune parameters, switch operating modes, select a different policy, or revert to a default configuration.

Both examples above (\Cref{lst:benefit-correction,lst:in-dist-short}) illustrate the {\sf Set} action: the CBMM guardrail repeatedly invokes {\sf Set} for estimated benefits of a page, and the in-distribution check guardrail sets the {\tt flux\_policy} flag to {\tt FLUXLongPredictor}.
The assumption for both of these examples is that setting these flags is consumed by the system correctly, and accordingly changes to the policy are made.
Rather than encoding arbitrary recovery logic in the DSL, \lang uses these actions as a narrow interface to policy-specific mechanisms already implemented in the system.

\paragraph{Verification Conditions.} \aditya{This is a little opaque and some new terms are introduced here? e.g., is the property in check the global objective or is it the one in ensures... can you label the example?}
\lang allows specification of optional verification conditions using the {\tt ensures} and {\tt assert} keywords, as shown in \Cref{lst:in-dist-short}.
When specified, the {\tt ensures} condition is first checked to hold whenever the property specified in {\tt check} holds, and whenever the property specified in {\tt check} does not hold but the action satisfies the condition.
Similarly, the {\tt assert} condition is verified to hold when all guardrails meet their respective {\tt ensures} conditions.
Both  {\tt ensures} and {\tt assert} statements are specified as uninterpreted functions in \lang, that must be given meaning via the system-policy interaction model ($M$ in \Cref{subsec:verification-conditions}).
In this example, both the local guarantee and the \perfsafety condition happen to be the same, but as we show in later examples, this may not always hold. \aditya{this somewhat clarifies my confusion, but the writing can be more direct}
